\section{Tone Deaf to Life}

\begin{quotex}
Man's life is a becoming; and not only becoming, but self-creation. He does not grow under the direction and control of irresistible forces. The force that shapes him is his own will. All his life is an effort to attain to real human nature. But human nature, since man is at bottom spirit, is only exemplified in the absolute spirit of God. Hence man must shape himself in God's image, or he ceases to be even human and becomes diabolical. This self-creation must also be self-knowledge; not the self-knowledge of introspection, the examination of the self that is, but the knowledge of God, the self that is to be. Knowledge of God is the beginning, the center, and end, of human life.
\flright{\textsc{R. G. Collingwood}}

Creation emanates from the potentialities inherent in God's being, a being which radiates ``into'' the coming into visibility and intelligibility of all that surrounds us.
\flright{\textsc{Ibn Arabi}}

Life is the actualization of possibilities.
\flright{\textsc{Ulick Varange}}

The worst thing that can happen to a man who is already divided up into a dozen different compartments is to seal off yet another compartment and tell him that this one is more important than all others, and that he must henceforth exercise a special care in keeping it separate from them.

The first thing that you have to do, before you even start thinking about such a thing as contemplation, is to try to recover your basic natural unity, reintegrate your compartmentalized being into a coordinated and simple whole and learn to live as a unified human person. This means that you have to bring back together the fragments of your distracted existence so that when you say ``I'', there really is someone present to support the pronoun you have uttered. Reflect, sometimes, on the disquieting fact that most of your statements of opinions, tastes, deeds, desires, hopes, and fears are statements about someone who is not really present. When you say ``I think,'' it is often not you who think, but ``they'' --- it is the anonymous authority of the collectivity speaking through your mask. When you say ``I want,'' you are accepting, paying for, what has been forced upon you. That is to say, you reach out for what you have been made to want.
\flright{\textsc{Thomas Merton}}
\end{quotex}


\paragraph{Symphony}


In studying creativity, we have used the example of the symphony as an aid to understand the Noumenon and its relationship to phenomena. We recognize a symphony by its unity ... notes follow each other in melody, transitions are resolved harmoniously, movements relate to each other, all within an overarching rhythm. Great composers have often described the creative process. They don't begin with the parts and end up with a unity. Quite the contrary ... they grasp the symphony in its wholeness and scramble to gets its parts down on paper.

For example, this is Mozart's description of the creative process:

\begin{quotex}
All this fires my soul and, provided I am not disturbed, my subject enlarges itself, becomes methodized and defined, and the whole, though it be long, stands almost complete and finished in my mind, so I can survey it, like a fine picture or a beautiful statue at a glance.
... All this inventing, this producing, takes place in a pleasing, lively dream.
\end{quotex}


Tchaikovsky offers a similar description:

\begin{quotex}
Generally speaking, the germ of a future composition comes suddenly and unexpectedly ... it takes root with extraordinary force and rapidity, shoots up through the earth, puts forth branches and leaves and finally blossoms. I cannot define the creative process in any other way than by this simile. ... I forget everything and behave like a madman: everything within me stands pulsing and quivering; hardly have I begun the sketch before one thought follows another. In the midst of this magic process, it frequently happens that some external interruption awakes me from my somnambulistic state. ... such dreadful interruptions break the thread of inspiration.
\end{quotex}


For us, the listeners, on the other hand, the process works in reverse. We hear the sounds: the instruments coming in and out, the melody line, the harmony, the percussion, the changes in volume and pitch. The amazing thing is that we don't focus on each individual sound, but rather we grasp the symphony as a whole -- usually it only takes a few notes and we recognize it immediately. ``Beethoven's 9th'', we call out.

But exactly where is the 9th symphony? Until it is performed, it has only a possible existence and it exists nowhere and everywhere. Once it is performed, it suddenly comes alive, it manifests itself, the possible becomes actual. What are the factors? A complex mix of necessity, fortune, freedom. The instruments are required. The mysterious element of fortune comes into play: why is one successful, the other not. The individual decisions of the players must converge. All that --- plus the symphony itself --- are part of the stew.

\paragraph{Life}


How do we experience our own lives? The traditional view of man is spirit, soul, and body. Where is the spirit? ``Before I formed you in the womb, I knew you, and before you were born I consecrated you'' (Jer 1:5) God knows us as possibility, as spirit. Traditionally, the soul is the form of the body. The soul, then, is like the notes of the symphony that is spirit, written down by the composer to match perfectly the unity of the symphony. The body, then, is the manifestation of spirit and soul, its coming into awareness.

Again, there are the same mysterious factors. Our destiny as determined by God. Fortune, the force of circumstances, the forces that shape and form our environment, that close out some options while opening up other possibilities. Then freedom, the choices we make. All this combines to make our life, to actualize the possibilities available to us. Is that how we see it? There are all the events in the day, from the trivial to the momentous. We forget what happened a year ago, we see no relation between today and a year from today. Our life is in fragments: a job, a family, friends, entertainment, education, etc. Do we pause to see how all the pieces of our lives fit together into a harmonious whole, how we are creating it year by year, day by day, minute by minute?

There is a disability called ``tone deafness''. A tone deaf person hears each note of the symphony, every sound, all the instruments; yet, he doesn't hear the symphony itself. To him, the notes are unrelated, merely a cacophony of individual sounds, little different from street noise.

In an analogous way, most people are \emph{Tone Deaf to Life}.


\flrightit{Posted on 2007-08-24 by Cologero}

